% Generated mechanically from frozen input inputs/p12/ja/stacks-limits.tex.
% Do not edit this generated file; edit the bound locale source instead.

% OK, start here.
%

\setcounter{chapter}{101}
\chapter{代数スタックの極限}



\phantomsection
\label{stacks-limits-section-phantom}


\section{はじめに}
\label{stacks-limits-section-introduction}

\noindent
本章では，代数スタックの極限に関する事項をまとめる．
代数スタックおよび代数空間の極限に関する多くの結果は，
David Rydh により \cite{rydh_approx} で得られている．


\section{規約}
\label{stacks-limits-section-conventions}

\noindent
引き続き，Properties of Stacks, Section
\ref{stacks-properties-section-conventions} で導入した規約および
言葉の濫用を用いる．







\section{有限表示の射}
\label{stacks-limits-section-finite-presentation}

\noindent
本節は Limits of Spaces, Section
\ref{spaces-limits-section-finite-presentation} の類似である．
そこでは，$\Sch$ 上の関手の変換が極限を保つとは何かを定義した
（Limits of Spaces, Lemma
\ref{spaces-limits-lemma-characterize-relative-limit-preserving}
における特徴づけを参照されたい）．
Criteria for Representability, Section
\ref{criteria-section-limit-preserving} では，
「対象について極限を保つ」という概念を定義した．
また Artin's Axioms, Section \ref{artin-section-limits} では，
$\Sch$ 上の亜群にファイバー化された圏が極限を保つとは何かを定義した．
これらを組み合わせると，次の概念を得る．

\begin{definition}
\label{stacks-limits-definition-limit-preserving}
$S$ をスキームとする．
$f : \mathcal{X} \to \mathcal{Y}$ を $(\Sch/S)_{fppf}$ 上の
亜群にファイバー化された圏の $1$-射とする．
$S$ 上のアフィンスキームの任意の有向極限
$U = \lim U_i$ に対し，ファイバー圏の図式
$$
\xymatrix{
\colim \mathcal{X}_{U_i} \ar[r] \ar[d]_f & \mathcal{X}_U \ar[d]^f \\
\colim \mathcal{Y}_{U_i} \ar[r] & \mathcal{Y}_U
}
$$
が $2$-カルテジアンであるとき，$f$ は {\it 極限を保つ} という．
\end{definition}

\begin{lemma}
\label{stacks-limits-lemma-limit-preserving-objects}
$S$ をスキームとする．
$f : \mathcal{X} \to \mathcal{Y}$ を $(\Sch/S)_{fppf}$ 上の
亜群にファイバー化された圏の $1$-射とする．
$f$ が極限を保つならば
（Definition \ref{stacks-limits-definition-limit-preserving}），
$f$ は対象について極限を保つ
（Criteria for Representability, Section
\ref{criteria-section-limit-preserving}）．
\end{lemma}

\begin{proof}
$U$ 上のアフィンスキームの任意の有向極限 $U = \lim U_i$ に対して，
関手
$$
\colim \mathcal{X}_{U_i} \longrightarrow
(\colim \mathcal{Y}_{U_i}) \times_{\mathcal{Y}_U} \mathcal{X}_U
$$
が本質的全射ならば，$f$ は対象について極限を保つ．
\end{proof}

\begin{lemma}
\label{stacks-limits-lemma-base-change-limit-preserving}
$p : \mathcal{X} \to \mathcal{Y}$ および
$q : \mathcal{Z} \to \mathcal{Y}$ を $(\Sch/S)_{fppf}$ 上の
亜群にファイバー化された圏の $1$-射とする．
$p : \mathcal{X} \to \mathcal{Y}$ が極限を保つならば，
$q$ による $p$ の基底変換
$p' : \mathcal{X} \times_\mathcal{Y} \mathcal{Z} \to \mathcal{Z}$
も極限を保つ．
\end{lemma}

\begin{proof}
これは形式的である．
$U = \lim_{i \in I} U_i$ を $S$ 上のアフィンスキーム $U_i$ の
有向極限とする．各 $i$ に対して
$$
(\mathcal{X} \times_\mathcal{Y} \mathcal{Z})_{U_i} =
\mathcal{X}_{U_i} \times_{\mathcal{Y}_{U_i}} \mathcal{Z}_{U_i}
$$
である．フィルター余極限は圏の $2$-ファイバー積と可換である
（詳細は省略する）．したがって，$p$ が極限を保つならば
\begin{align*}
\colim (\mathcal{X} \times_\mathcal{Y} \mathcal{Z})_{U_i}
& =
\colim \mathcal{X}_{U_i} \times_{\colim \mathcal{Y}_{U_i}}
\colim \mathcal{Z}_{U_i} \\
& =
\mathcal{X}_U \times_{\mathcal{Y}_U} \colim \mathcal{Y}_{U_i}
\times_{\colim \mathcal{Y}_{U_i}}
\colim \mathcal{Z}_{U_i} \\
& =
\mathcal{X}_U \times_{\mathcal{Y}_U} \colim \mathcal{Z}_{U_i} \\
& =
\mathcal{X}_U \times_{\mathcal{Y}_U} \mathcal{Z}_U \times_{\mathcal{Z}_U}
\colim \mathcal{Z}_{U_i} \\
& =
(\mathcal{X} \times_\mathcal{Y} \mathcal{Z})_U \times_{\mathcal{Z}_U}
\colim \mathcal{Z}_{U_i}
\end{align*}
となり，所望の結論を得る．
\end{proof}

\begin{lemma}
\label{stacks-limits-lemma-composition-limit-preserving}
$p : \mathcal{X} \to \mathcal{Y}$ および
$q : \mathcal{Y} \to \mathcal{Z}$ を $(\Sch/S)_{fppf}$ 上の
亜群にファイバー化された圏の $1$-射とする．
$p$ と $q$ が極限を保つならば，合成 $q \circ p$ も極限を保つ．
\end{lemma}

\begin{proof}
これは形式的である．
$U = \lim_{i \in I} U_i$ を $S$ 上のアフィンスキーム $U_i$ の
有向極限とする．$p$ と $q$ が極限を保つならば
\begin{align*}
\colim \mathcal{X}_{U_i}
& =
\mathcal{X}_U \times_{\mathcal{Y}_U} \colim \mathcal{Y}_{U_i} \\
& =
\mathcal{X}_U \times_{\mathcal{Y}_U} \mathcal{Y}_U
\times_{\mathcal{Z}_U} \colim \mathcal{Z}_{U_i} \\
& =
\mathcal{X}_U \times_{\mathcal{Z}_U} \colim \mathcal{Z}_{U_i}
\end{align*}
となり，所望の結論を得る．
\end{proof}

\begin{lemma}
\label{stacks-limits-lemma-representable-by-spaces-limit-preserving}
$p : \mathcal{X} \to \mathcal{Y}$ を $(\Sch/S)_{fppf}$ 上の
亜群にファイバー化された圏の $1$-射とする．
$p$ が代数空間によって表現可能ならば，次の条件は同値である：
\begin{enumerate}
\item $p$ は極限を保つ，
\item $p$ は対象について極限を保つ，
\item $p$ は局所有限表示である
（Algebraic Stacks,
Definition \ref{algebraic-definition-relative-representable-property} を参照）．
\end{enumerate}
\end{lemma}

\begin{proof}
Criteria for Representability, Lemma
\ref{criteria-lemma-representable-by-spaces-limit-preserving} において，
(2) と (3) が同値であることを見た．
したがって，(1) と (2) が同値であることを示せば十分である．
一方の向きは Lemma \ref{stacks-limits-lemma-limit-preserving-objects} で見た．
逆向きについて，$U = \lim_{i \in I} U_i$ を
$S$ 上のアフィンスキーム $U_i$ の有向極限とする．
示すべきことは
$$
\colim \mathcal{X}_{U_i} \longrightarrow
\mathcal{X}_U \times_{\mathcal{Y}_U} \colim \mathcal{Y}_{U_i}
$$
が同値であることである．(2) を仮定しているので，
これが本質的全射であることは分かっている．
したがって，充満忠実であることを証明すればよい．
$p$ はファイバー圏上で忠実である
（Algebraic Stacks, Lemma
\ref{algebraic-lemma-criterion-map-representable-spaces-fibred-in-groupoids}）
から，上の関手は忠実である．
$x_i$ および $x'_i$ を $U_i$ 上の $\mathcal{X}$ の
ファイバー圏の対象とする．
上の関手は $x_i$ を $(x_i|_U, p(x_i), can)$ に送る．
ここで $can$ は標準同型 $p(x_i|_U) \to p(x_i)|_U$ である．
そこで，この証明で最初に表示した矢印の右辺の圏における射
$$
(\alpha, \beta_i) : (x_i|_U, p(x_i), can) \longrightarrow
(x'_i|_U, p(x'_i), can)
$$
が与えられていると仮定する．
われわれの課題は，$i' \geq i$ と，$(\alpha, \beta_i|_{U_{i'}})$ に
写る射 $x_i|_{U_{i'}} \to x'_i|_{U_{i'}}$ を構成することである．

\medskip\noindent
$y_i = p(x_i)$ および $y'_i = p(x'_i)$ とおく．
Algebraic Stacks, Lemma
\ref{algebraic-lemma-criterion-map-representable-spaces-fibred-in-groupoids}
により，関手
$$
X_{y_i} : (\Sch/U_i)^{opp} \to \textit{Sets},\quad
V/U_i \mapsto
\{(x, \phi) \mid x \in \Ob(\mathcal{X}_V), \phi : f(x) \to y_i|V\}/\cong
$$
は $U_i$ 上の代数空間であり，同様に定義した関手
$X_{y'_i}$ についても同じことが成り立つ．
(2) と (3) は同値なので，$X_{y'_i}$ は $U_i$ 上局所有限表示である．
さらに，$(x_i, \text{id})$ および $(x'_i, \text{id})$ はそれぞれ
$X_{y_i}$ および $X_{y'_i}$ の $U_i$-値点を定める．
関手の変換
$$
\beta_i : X_{y_i} \to X_{y'_i},\quad
(x/V, \phi) \mapsto (x/V, \beta_i|_V \circ \phi)
$$
が存在する．言い換えれば，これは $U_i$ 上の代数空間の射である．
次の図式は可換であると主張する：
$$
\xymatrix{
U \ar[d] \ar[rr] & & U_i \ar[d]^{(x'_i, \text{id})} \\
U_i \ar[r]^{(x_i, \text{id})} & X_{y_i} \ar[r]^{\beta_i} & X_{y'_i}
}
$$
実際，これは，関手 $X_{y'_i}$ の定義に現れる対
$(x_i|_U, \beta_i|_U)$ と $(x'_i|_U, \text{id})$ が
同型であるという条件に同値である．
そして射 $\alpha : x_i|_U \to x'_i|_U$ が，まさにそのような同型を与える．
逆にたどれば，次の図式が可換となるような $i' \geq i$ を見つけられれば，
$$
\xymatrix{
U_{i'} \ar[d] \ar[rr] & & U_i \ar[d]^{(x'_i, \text{id})} \\
U_i \ar[r]^{(x_i, \text{id})} & X_{y_i} \ar[r]^{\beta_i} & X_{y'_i}
}
$$
前段落で提示した問題の解となる同型
$x_i|_{U_{i'}} \to x'_i|_{U_{i'}}$ が得られることが分かる．
ところが，対角射
$$
\Delta : X_{y'_i} \to X_{y'_i} \times_{U_i} X_{y'_i}
$$
は局所有限表示である
（Morphisms of Spaces, Lemma
\ref{spaces-morphisms-lemma-diagonal-morphism-finite-type}）．
したがって，$U \to U_i$ が $X_{y'_i}$ への二つの射を等化するという事実から，
ある $i' \geq i$ に対して $U_{i'} \to U_i$ がそれらを等化する．
Limits of Spaces, Proposition
\ref{spaces-limits-proposition-characterize-locally-finite-presentation}
を参照されたい．
\end{proof}

\begin{lemma}
\label{stacks-limits-lemma-limit-preserving-diagonal}
$p : \mathcal{X} \to \mathcal{Y}$ を $(\Sch/S)_{fppf}$ 上の
亜群にファイバー化された圏の $1$-射とする．
次の条件は同値である：
\begin{enumerate}
\item 対角射
$\Delta : \mathcal{X} \to \mathcal{X} \times_\mathcal{Y} \mathcal{X}$
は極限を保つ，
\item $S$ 上のアフィンスキームの任意の有向極限
$U = \lim U_i$ に対して，関手
$$
\colim \mathcal{X}_{U_i} \longrightarrow
\mathcal{X}_U \times_{\mathcal{Y}_U} \colim \mathcal{Y}_{U_i}
$$
は充満忠実である．
\end{enumerate}
特に，$p$ が極限を保つならば，$\Delta$ も極限を保つ．
\end{lemma}

\begin{proof}
$U = \lim U_i$ を $S$ 上のアフィンスキームの有向極限とする．
関手
$$
\colim \mathcal{X}_{U_i} \longrightarrow
\mathcal{X}_U \times_{\mathcal{Y}_U} \colim \mathcal{Y}_{U_i}
$$
が充満忠実であることと，関手
$$
\colim \mathcal{X}_{U_i} \longrightarrow
\mathcal{X}_U \times_{(\mathcal{X} \times_\mathcal{Y} \mathcal{X})_U}
\colim (\mathcal{X} \times_\mathcal{Y} \mathcal{X})_{U_i}
$$
が同値であることは同値であると主張する．
これを示せば補題が従う．
実際，
$(\mathcal{X} \times_\mathcal{Y} \mathcal{X})_U =
\mathcal{X}_U \times_{\mathcal{Y}_U} \mathcal{X}_U$
かつ
$(\mathcal{X} \times_\mathcal{Y} \mathcal{X})_{U_i} =
\mathcal{X}_{U_i} \times_{\mathcal{Y}_{U_i}} \mathcal{X}_{U_i}$
であるから，これは純粋に圏論的な主張である．
次の段落でこれを論じる．

\medskip\noindent
$\mathcal{I}$ をフィルター付き添字圏とする．
$(\mathcal{C}_i)$ および $(\mathcal{D}_i)$ を
$\mathcal{I}$ 上の亜群の系とし，
$p : (\mathcal{C}_i) \to (\mathcal{D}_i)$ を
$\mathcal{I}$ 上の亜群の系の射とする．
さらに亜群の関手 $p : \mathcal{C} \to \mathcal{D}$ と，
可換図式
$$
\xymatrix{
\colim \mathcal{C}_i \ar[d]_p \ar[r]_f & \mathcal{C} \ar[d]^p \\
\colim \mathcal{D}_i \ar[r]^g & \mathcal{D}
}
$$
をなす関手
$f : \colim \mathcal{C}_i \to \mathcal{C}$ および
$g : \colim \mathcal{D}_i \to \mathcal{D}$ が与えられているとする．
このとき
$$
A : \colim \mathcal{C}_i \longrightarrow
\mathcal{C} \times_\mathcal{D} \colim \mathcal{D}_i
$$
が充満忠実であることと，関手
$$
B : \colim \mathcal{C}_i \longrightarrow
\mathcal{C}
\times_{\Delta, \mathcal{C} \times_\mathcal{D} \mathcal{C}, f \times_g f}
\colim (\mathcal{C}_i \times_{\mathcal{D}_i} \mathcal{C}_i)
$$
が同値であることは同値であると主張する．
$\mathcal{C}' = \colim \mathcal{C}_i$ および
$\mathcal{D}' = \colim \mathcal{D}_i$ とおく．
$2$-ファイバー積はフィルター余極限と可換するので，
$A$ と $B$ はそれぞれ関手
$$
A' : \mathcal{C}' \to \mathcal{C} \times_\mathcal{D} \mathcal{D}'
\quad\text{and}\quad
B' : \mathcal{C}' \longrightarrow
\mathcal{C}
\times_{\Delta, \mathcal{C} \times_\mathcal{D} \mathcal{C}, f \times_g f}
(\mathcal{C}' \times_{\mathcal{D}'} \mathcal{C}')
$$
となる．したがって，
$$
\xymatrix{
\mathcal{C}' \ar[d]_p \ar[r]_f & \mathcal{C} \ar[d]^p \\
\mathcal{D}' \ar[r]^g & \mathcal{D}
}
$$
が亜群の可換図式ならば，$A'$ が充満忠実であることと
$B'$ が同値であることが同値であると示せば十分である．
これは Categories, Lemma
\ref{categories-lemma-fully-faithful-diagonal-equivalence}
（基底圏を自明，すなわち一点からなる圏とする）から従う．実際，
$$
\mathcal{C}
\times_{\Delta, \mathcal{C} \times_\mathcal{D} \mathcal{C}, f \times_g f}
(\mathcal{C}' \times_{\mathcal{D}'} \mathcal{C}') =
\mathcal{C}'
\times_{A', \mathcal{C} \times_\mathcal{D} \mathcal{D}', A'}
\mathcal{C}'
$$
である．これで証明が完了する．
\end{proof}

\begin{lemma}
\label{stacks-limits-lemma-locally-finite-presentation-limit-preserving}
$S$ をスキームとし，$\mathcal{X}$ を $S$ 上の代数スタックとする．
$\mathcal{X} \to S$ が局所有限表示ならば，$\mathcal{X}$ は
Artin's Axioms, Definition \ref{artin-definition-limit-preserving}
の意味で極限を保つ
（同値な言い方をすれば，射 $\mathcal{X} \to S$ は極限を保つ）．
\end{lemma}

\begin{proof}
あるスキーム $U$ からの全射な滑らかな射
$U \to \mathcal{X}$ を選ぶ．このとき $U \to S$ は局所有限表示である．
Morphisms of Stacks, Section
\ref{stacks-morphisms-section-finite-presentation} を参照されたい．
代数空間のある滑らかな亜群 $(U, R, s, t, c)$ により
$\mathcal{X} = [U/R]$ と書ける．
Algebraic Stacks, Lemma \ref{algebraic-lemma-stack-presentation}
を参照されたい．$U$ は $S$ 上局所有限表示なので，
代数空間 $R$ も $S$ 上局所有限表示である．
$[U/R]$ は，$T$ 上のファイバー圏が亜群
$(U(T), R(T), s, t, c)$ であるような亜群にファイバー化された圏を
スタック化して得られる，$(\Sch/S)_{fppf}$ 上の亜群のスタックであることを
思い出そう．$U$ と $R$ は関手として極限を保つので
（Limits of Spaces, Proposition
\ref{spaces-limits-proposition-characterize-locally-finite-presentation}），
この亜群にファイバー化された圏は極限を保つ．
したがって，fppf スタック化が極限を保つという性質を保存することを
示せば十分である．これは正しい
（ヒント：Topologies, Lemma
\ref{topologies-lemma-limit-fppf-topology} を用いよ）．
しかし以下では，この場合にスタック化が何を意味するかが分かっていることを
用いて，直接の証明を与える．

\medskip\noindent
$T = \lim T_\lambda$ を $S$ 上のアフィンスキームの有向極限とする．
関手
$$
\colim [U/R]_{T_\lambda} \longrightarrow [U/R]_T
$$
が圏の同値であることを示さなければならない．
まずこの関手が本質的全射であることを示す．
$x \in \Ob([U/R]_T)$ とする．
Groupoids in Spaces, Lemma
\ref{spaces-groupoids-lemma-quotient-stack-objects} に，
圏 $[U/R]_T$ の記述がある．
特に，$x$ は fppf 被覆 $\{T_i \to T\}_{i \in I}$ と，
この被覆に関する $[U/R]$-降下データ $(u_i, r_{ij})$ に対応する．
この被覆を細分することにより，アフィンスキーム $T$ の
標準 fppf 被覆であると仮定してよい．
Topologies, Lemma
\ref{topologies-lemma-limit-fppf-topology} により，
ある $\lambda$ と標準 fppf 被覆
$\{T_{\lambda, i} \to T_\lambda\}_{i \in I}$ を選んで，
$T$ への基底変換が $\{T_i \to T\}_{i \in I}$ に等しくなるようにできる．
各 $i$ に対し，$\lambda$ を増大させれば，
$T_i \to T_{\lambda, i}$ との合成が与えられた射 $u_i$ となる
$u_{\lambda, i} : T_{\lambda, i} \to U$ を見つけられる
（ここで $U$ が極限を保つことを用いる）．
同様に，各 $i, j$ に対し，$\lambda$ を増大させれば，
$T_{ij} \to T_{\lambda, ij}$ との合成が与えられた射 $r_{ij}$ となる
$r_{\lambda, ij} : T_{\lambda, i} \times_{T_\lambda} T_{\lambda, j} \to R$
を見つけられる（ここで $R$ が極限を保つことを用いる）．
さらに $\lambda$ を増大させて，
$$
s \circ r_{\lambda, ij} = u_{\lambda, i} \circ \text{pr}_0
\quad\text{and}\quad
t \circ r_{\lambda, ij} = u_{\lambda, j} \circ \text{pr}_1,
$$
および
$$
c \circ (r_{\lambda, jk} \circ \text{pr}_{12}, r_{\lambda, ij}
\circ \text{pr}_{01}) = r_{\lambda, ik} \circ \text{pr}_{02}.
$$
を仮定してよい．言い換えれば，$(u_{\lambda, i}, r_{\lambda, ij})$ が
被覆 $\{T_{\lambda, i} \to T_\lambda\}_{i \in I}$ に関する
$[U/R]$-降下データであると仮定してよい．
すると $T_\lambda$ 上の $[U/R]$ の対応する対象が得られ，
その $T$ への引き戻しは所望どおり $x$ と同型になる．
充満忠実性の証明も，Groupoids in Spaces, Lemma
\ref{spaces-groupoids-lemma-quotient-stack-objects} に与えられた
$[U/T]$ のファイバー圏における射の記述を用いて，まったく同様に進む．
\end{proof}

\begin{proposition}
\label{stacks-limits-proposition-characterize-locally-finite-presentation}
\begin{reference}
これは \cite[Lemma 2.3.15]{Emerton-Gee} の特別な場合である．
\end{reference}
代数スタックの射 $f : \mathcal{X} \to \mathcal{Y}$ に対し，
次の条件は同値である：
\begin{enumerate}
\item $f$ は極限を保つ，
\item $f$ は対象について極限を保つ，
\item $f$ は局所有限表示である．
\end{enumerate}
\end{proposition}

\begin{proof}
(3) を仮定する．
$T = \lim T_i$ をアフィンスキームの有向極限とし，関手
$$
\colim \mathcal{X}_{T_i} \longrightarrow
\mathcal{X}_T \times_{\mathcal{Y}_T} \colim \mathcal{Y}_{T_i}
$$
を考える．右辺の対象 $(x, y_i, \beta)$，すなわち
$x \in \Ob(\mathcal{X}_T)$，$y_i \in \Ob(\mathcal{Y}_{T_i})$，および
$\mathcal{Y}_T$ における $\beta : f(x) \to y_i|_T$ を考える．
$(x, y_i, \beta)$ は，$T$ 上の代数スタック
$\mathcal{X}_{y_i} = \mathcal{X} \times_{\mathcal{Y}, y_i} T_i$
の対象とみなせる．
$\mathcal{X}_{y_i} \to T_i$ は $f$ の基底変換として局所有限表示なので，
Lemma \ref{stacks-limits-lemma-locally-finite-presentation-limit-preserving} により
極限を保つ．これは，ある $i' \geq i$ に対して
$(x, y_i, \beta)$ が $T_{i'}$ 上の対象から来ることを意味する．
定義を展開すれば，$(x, y_i, \beta)$ は表示された関手の
本質像に属することが分かる．
言い換えれば，その関手は本質的全射であり，
これは $f$ が対象について極限を保つということでもある．
次に，これを $f$ の対角射 $\Delta$ に適用する．
Morphisms of Stacks, Lemma
\ref{stacks-morphisms-lemma-diagonal-morphism-finite-type} により，
射 $\Delta$ は局所有限表示である．
したがって，上の議論から $\Delta$ は対象について極限を保つ．
Lemma \ref{stacks-limits-lemma-representable-by-spaces-limit-preserving} により，
$\Delta$ は極限を保つ．さらに
Lemma \ref{stacks-limits-lemma-limit-preserving-diagonal} により，
上で表示した関手は充満忠実である．
すでに本質的全射であることを示したので，これは同値であり，
(1) が成り立つ．

\medskip\noindent
(1) $\Rightarrow$ (2) は自明である．(2) を仮定する．
スキーム $V$ と全射な滑らかな射 $V \to \mathcal{Y}$ を選ぶ．
Criteria for Representability, Lemma
\ref{criteria-lemma-base-change-limit-preserving} により，
基底変換 $\mathcal{X} \times_\mathcal{Y} V \to V$ は
対象について極限を保つ．
スキーム $U$ と全射な滑らかな射
$U \to \mathcal{X} \times_\mathcal{Y} V$ を選ぶ．
滑らかな射は局所有限表示なので，
$U \to \mathcal{X} \times_\mathcal{Y} V$ は極限を保つ
（証明の前半）．
Criteria for Representability, Lemma
\ref{criteria-lemma-composition-limit-preserving} により，
合成 $U \to V$ は対象について極限を保つ．
Criteria for Representability, Lemma
\ref{criteria-lemma-representable-by-spaces-limit-preserving} により，
$U \to V$ は局所有限表示である．
これは，まさに $f$ が局所有限表示であるための条件である．
Morphisms of Stacks, Definition
\ref{stacks-morphisms-definition-locally-finite-presentation}
を参照されたい．
\end{proof}




\section{性質の降下}
\label{stacks-limits-section-descent}

\noindent
本節は Limits, Section \ref{limits-section-descent} の類似である．

\begin{situation}
\label{stacks-limits-situation-descent}
$Y = \lim_{i \in I} Y_i$ を，遷移射がアフィンである
代数空間の有向系の極限とする．
すべての $i \in I$ に対して $X_i$ は準コンパクトかつ準分離であると仮定する．
さらに元 $0 \in I$ を一つ選ぶ．
\end{situation}

\begin{lemma}
\label{stacks-limits-lemma-eventually-separated}
Situation \ref{stacks-limits-situation-descent} において，
$\mathcal{X}_0 \to Y_0$ を代数スタックから $Y_0$ への射とする．
$\mathcal{X}_0$ は準コンパクトかつ準分離であると仮定する．
$Y \times_{Y_0} \mathcal{X}_0 \to Y$ が分離的ならば，
十分大きいすべての $i \in I$ に対して
$Y_i \times_{Y_0} \mathcal{X}_0 \to Y_i$ は分離的である．
\end{lemma}

\begin{proof}
$\mathcal{X} = Y \times_{Y_0} \mathcal{X}_0$ および
$\mathcal{X}_i = Y_i \times_{Y_0} \mathcal{X}_0$ とおく．
アフィンスキーム $U_0$ と全射な滑らかな射
$U_0 \to \mathcal{X}_0$ を選ぶ．
$U = Y \times_{Y_0} U_0$ および
$U_i = Y_i \times_{Y_0} U_0$ とおく．
すると $U$ と $U_i$ はアフィンであり，
$U \to \mathcal{X}$ と $U_i \to \mathcal{X}_i$ は滑らかかつ全射である．
$R_0 = U_0 \times_{\mathcal{X}_0} U_0$ とおき，
$R = Y \times_{Y_0} R_0$ および $R_i = Y_i \times_{Y_0} R_0$ とおく．
このとき $R = U \times_\mathcal{X} U$ かつ
$R_i = U_i \times_{\mathcal{X}_i} U_i$ である．

\medskip\noindent
この記法のもとで，$\mathcal{X} \to Y$ が分離的ならば，
$R \to U \times_Y U$ は
$U \times_Y U \to \mathcal{X} \times_Y \mathcal{X}$ による
$\mathcal{X} \to \mathcal{X} \times_Y \mathcal{X}$ の基底変換として
固有である．逆に，
$U_i \times_{Y_i} U_i \to \mathcal{X}_i \times_{Y_i} \mathcal{X}_i$
は全射かつ滑らかなので，
$R_i \to U_i \times_{Y_i} U_i$ が固有ならば
$\mathcal{X}_i \to Y_i$ は分離的である．
Properties of Stacks, Lemma
\ref{stacks-properties-lemma-check-property-covering} を参照されたい．
$R_0 \to U_0 \times_{Y_0} U_0$ は局所有限型であり，
$R_0$ は準コンパクトかつ準分離であることに注意する．
Limits of Spaces, Lemma \ref{spaces-limits-lemma-eventually-proper}
により，十分大きい $i$ に対して
$R_i \to U_i \times_{Y_i} U_i$ は固有である．
これで証明が完了する．
\end{proof}







\section{相対的対象の降下}
\label{stacks-limits-section-descending-relative}

\noindent
本節は Limits of Spaces, Section
\ref{spaces-limits-section-descending-relative} の類似である．

\begin{lemma}
\label{stacks-limits-lemma-descend-a-stack-down}
$I$ を有向集合とし，$(X_i, f_{ii'})$ を $I$ 上の代数空間の逆系とする．
次を仮定する：
\begin{enumerate}
\item 射 $f_{ii'} : X_i \to X_{i'}$ はアフィンである，
\item 空間 $X_i$ は準コンパクトかつ準分離である．
\end{enumerate}
$X = \lim X_i$ とする．
$\mathcal{X}$ が $X$ 上有限表示の代数スタックならば，
ある $i \in I$ と $X_i$ 上有限表示の代数スタック $\mathcal{X}_i$ が存在し，
$X$ 上の代数スタックとして
$\mathcal{X} \cong \mathcal{X}_i \times_{X_i} X$ となる．
\end{lemma}

\begin{proof}
Morphisms of Stacks, Definition
\ref{stacks-morphisms-definition-locally-finite-presentation} により，
射 $\mathcal{X} \to X$ は準コンパクト，局所有限表示，かつ準分離である．
$X$ は準コンパクトで $\mathcal{X} \to X$ は準コンパクトなので，
$\mathcal{X}$ は準コンパクトである
（Morphisms of Stacks, Definition
\ref{stacks-morphisms-definition-quasi-compact}）．
したがって，アフィンスキーム $U$ と全射な滑らかな射
$U \to \mathcal{X}$ が存在する
（Properties of Stacks, Lemma
\ref{stacks-properties-lemma-quasi-compact-stack}）．
$R = U \times_\mathcal{X} U$ とおく．
すると $X$ 上の代数空間の滑らかな亜群 $(U, R, s, t, c)$ で
$\mathcal{X} = [U/R]$ となるものを得る．
Algebraic Stacks, Lemma
\ref{algebraic-lemma-stack-presentation} を参照されたい．
$\mathcal{X} \to X$ は準分離で $X$ も準分離なので，
$\mathcal{X}$ は準分離である
（Morphisms of Stacks, Lemma
\ref{stacks-morphisms-lemma-composition-separated}）．
したがって $R \to U \times U$ は準コンパクトかつ準分離であり
（Morphisms of Stacks, Lemma
\ref{stacks-morphisms-lemma-fibre-product-after-map}），
$R$ は準分離かつ準コンパクトな代数空間である．
一方，$U \to X$ は局所有限表示であり，
$s : R \to U$ は滑らか，したがって局所有限表示なので，
$R \to X$ も局所有限表示である．
よって $(U, R, s, t, c)$ は，
$X$ 上有限表示である代数空間の圏における亜群対象である．
Limits of Spaces, Lemma
\ref{spaces-limits-lemma-descend-finite-presentation} により，
ある $i$ と，$X_i$ 上の代数空間の亜群
$(U_i, R_i, s_i, t_i, c_i)$ が存在し，
その $X$ への引き戻しは $(U, R, s, t, c)$ と同型になる．
$i$ を増大させれば，$s_i$ と $t_i$ は滑らかであると仮定できる．
Limits of Spaces, Lemma \ref{spaces-limits-lemma-descend-smooth}
を参照されたい．商スタック
$\mathcal{X}_i = [U_i/R_i]$ は代数スタックである
（Algebraic Stacks, Theorem
\ref{algebraic-theorem-smooth-groupoid-gives-algebraic-stack}）．

\medskip\noindent
射 $[U/R] \to [U_i/R_i]$ が存在する．
Groupoids in Spaces, Lemma
\ref{spaces-groupoids-lemma-quotient-stack-functorial} を参照されたい．
この射と射 $[U/R] \to X$ および $[U_i/R_i] \to X_i$
（Groupoids in Spaces, Lemma
\ref{spaces-groupoids-lemma-quotient-stack-arrows}）を合わせると，
同型（すなわち同値）
$$
[U/R] \longrightarrow [U_i/R_i] \times_{X_i} X
$$
が得られると主張する．
Groupoids in Spaces, Equation
(\ref{spaces-groupoids-equation-quotient-stack}) における
「亜群の前層」のレベルで，対応する射
$$
[U/_{\!p}R] \longrightarrow [U_i/_{\!p}R_i] \times_{X_i} X
$$
は同型である．したがって，スタック化がファイバー積と可換することから
主張が従う．Stacks, Lemma
\ref{stacks-lemma-stackification-fibre-product-fibred-categories}
を参照されたい．
\end{proof}




\section{有限表示の中への有限型の閉埋め込み}
\label{stacks-limits-section-finite-type-closed-in-finite-presentation}

\noindent
本節は Limits of Spaces, Section
\ref{spaces-limits-section-finite-type-closed-in-finite-presentation}
の類似である．

\begin{lemma}
\label{stacks-limits-lemma-finite-type-closed-in-finite-presentation}
$f : \mathcal{X} \to Y$ を代数スタックから代数空間への射とする．
次を仮定する：
\begin{enumerate}
\item $f$ は有限型かつ準分離である，
\item $Y$ は準コンパクトかつ準分離である．
\end{enumerate}
このとき，有限表示の射 $f' : \mathcal{X}' \to Y$ と，
$Y$ 上の代数スタックの閉埋め込み
$\mathcal{X} \to \mathcal{X}'$ が存在する．
\end{lemma}

\begin{proof}
この代数空間を，アフィンな遷移射をもつ有向集合 $I$ 上の代数空間の
極限 $Y = \lim_{i \in I} Y_i$ と書く．ここで $Y_i$ はネーターとする．
Limits of Spaces, Proposition
\ref{spaces-limits-proposition-approximate} を参照されたい．
Limits of Spaces, Section
\ref{spaces-limits-section-finite-type-quasi-separated}
の内容を用いる．

\medskip\noindent
表示 $\mathcal{X} = [U/R]$ を選ぶ．
対応する $Y$ 上の代数空間の亜群を
$(U, R, s, t, c, e, i)$ と表す．
$U$ はアフィンであると仮定してよく，実際そう仮定する．
このとき $U$，$R$，$R \times_{s, U, t} R$ は，
$Y$ 上有限型の準分離な代数空間である．
二つの射 $s, t : R \to U$，三つの射
$c : R \times_{s, U, t} R \to R$，
$\text{pr}_1 : R \times_{s, U, t} R \to R$，
$\text{pr}_2 : R \times_{s, U, t} R \to R$，
一つの射 $e : U \to R$，そして最後に一つの射 $i : R \to R$
がある．これらの射は Groupoids, Section
\ref{groupoids-section-groupoids} に詳述されている一連の公理を満たす．

\medskip\noindent
Limits of Spaces, Remark
\ref{spaces-limits-remark-finite-type-gives-well-defined-system} により，
ある $i_0 \in I$ と，$(Y_i)_{i \geq i_0}$ 上の逆系
\begin{enumerate}
\item $(U_i)_{i \geq i_0}$，
\item $(R_i)_{i \geq i_0}$，
\item $(T_i)_{i \geq i_0}$
\end{enumerate}
で，
$U = \lim_{i \geq i_0} U_i$，
$R = \lim_{i \geq i_0} R_i$，および
$R \times_{s, U, t} R = \lim_{i \geq i_0} T_i$
を満たすものを見つけられる．さらに系の射
\begin{enumerate}
\item $(s_i)_{i \geq i_0} : (R_i)_{i \geq i_0} \to (U_i)_{i \geq i_0}$，
\item $(t_i)_{i \geq i_0} : (R_i)_{i \geq i_0} \to (U_i)_{i \geq i_0}$，
\item $(c_i)_{i \geq i_0} : (T_i)_{i \geq i_0} \to (R_i)_{i \geq i_0}$，
\item $(p_i)_{i \geq i_0} : (T_i)_{i \geq i_0} \to (R_i)_{i \geq i_0}$，
\item $(q_i)_{i \geq i_0} : (T_i)_{i \geq i_0} \to (R_i)_{i \geq i_0}$，
\item $(e_i)_{i \geq i_0} : (U_i)_{i \geq i_0} \to (R_i)_{i \geq i_0}$，
\item $(i_i)_{i \geq i_0} : (R_i)_{i \geq i_0} \to (R_i)_{i \geq i_0}$
\end{enumerate}
が存在し，
$s = \lim_{i \geq i_0} s_i$，
$t = \lim_{i \geq i_0} t_i$，
$c = \lim_{i \geq i_0} c_i$，
$\text{pr}_1 = \lim_{i \geq i_0} p_i$，
$\text{pr}_2 = \lim_{i \geq i_0} q_i$，
$e = \lim_{i \geq i_0} e_i$，および
$i = \lim_{i \geq i_0} i_i$
となる．
Limits of Spaces, Lemma
\ref{spaces-limits-lemma-morphism-good-diagram-smooth} により，
$s_i$ と $t_i$ は滑らかであると仮定してよい
（そのために $i_0$ を増大させる必要があるかもしれない）．
Limits of Spaces, Lemma
\ref{spaces-limits-lemma-morphism-good-diagram-flat} により，
$s$ と $R \to R_i$ が定める射
$R \to U \times_{U_i, s_i} R_i$，および
$t$ と $R \to R_i$ が定める射
$R \to U \times_{U_i, t_i} R_i$ は，
すべての $i \geq i_0$ に対して同型であると仮定してよい．
Limits of Spaces, Lemma
\ref{spaces-limits-lemma-good-diagram-fibre-product} により，
図式
$$
\xymatrix{
T_i \ar[r]_{q_i} \ar[d]_{p_i} & R_i \ar[d]^{t_i} \\
R_i \ar[r]^{s_i} & U_i
}
$$
はカルテジアンであると仮定してよい．
すると Limits of Spaces, Lemma
\ref{spaces-limits-lemma-morphism-good-diagram} の一意性により，
十分大きい $i$ に対して，上で述べた射 $s, t, c, e, i$ の間の関係を
$s_i, t_i, c_i, e_i, i_i$ も満たすことが保証される．
そのような $i$ を一つ固定する．

\medskip\noindent
したがって，$(U_i, R_i, s_i, t_i, c_i, e_i, i_i)$ は
$Y_i$ 上の代数空間の滑らかな亜群である．
よって $\mathcal{X}_i = [U_i/R_i]$ は代数スタックである
（Algebraic Stacks, Theorem
\ref{algebraic-theorem-smooth-groupoid-gives-algebraic-stack}）．
$Y \to Y_i$ 上の亜群の射
$$
(U, R, s, t, c, e, i) \to (U_i, R_i, s_i, t_i, c_i, e_i, i_i)
$$
は可換図式
$$
\xymatrix{
\mathcal{X} \ar[d] \ar[r] & \mathcal{X}_i \ar[d] \\
Y \ar[r] & Y_i
}
$$
を定める
（Groupoids in Spaces, Lemma
\ref{spaces-groupoids-lemma-quotient-stack-functorial}）．
射 $\mathcal{X} \to Y \times_{Y_i} \mathcal{X}_i$ は
閉埋め込みであると主張する．
構成により代数スタック $\mathcal{X}_i \to Y_i$ は有限表示なので，
この主張から証明が完了する．
主張を示すため，左の図式
$$
\xymatrix{
U \ar[d] \ar[r] & U_i \ar[d] \\
\mathcal{X} \ar[r] & \mathcal{X}_i
}
\quad\quad
\xymatrix{
U \ar[d] \ar[r] & Y \times_{Y_i} U_i \ar[d] \\
\mathcal{X} \ar[r] & Y \times_{Y_i} \mathcal{X}_i
}
$$
は Groupoids in Spaces, Lemma
\ref{spaces-groupoids-lemma-criterion-fibre-product} と上の結果により
カルテジアンであることに注意する．
したがって右の可換図式もカルテジアンである．
ここで $U \to Y \times_{Y_i} U_i$ は
Limits of Spaces, Lemma
\ref{spaces-limits-lemma-limit-from-good-diagram} における逆系
$(U_i)$ の構成により閉埋め込みであり，
$Y \times_{Y_i} U_i \to Y \times_{Y_i} \mathcal{X}_i$ は
滑らかかつ全射である．
よって Properties of Stacks, Lemma
\ref{stacks-properties-lemma-check-immersion-covering} から
所望の結果が従う．
\end{proof}

\noindent
分離的な代数スタックに対する版もある．

\begin{lemma}
\label{stacks-limits-lemma-separated-closed-in-finite-presentation}
$f : \mathcal{X} \to Y$ を代数スタックから代数空間への射とする．
次を仮定する：
\begin{enumerate}
\item $f$ は有限型かつ分離的である，
\item $Y$ は準コンパクトかつ準分離である．
\end{enumerate}
このとき，有限表示の分離射 $f' : \mathcal{X}' \to Y$ と，
$Y$ 上の代数スタックの閉埋め込み
$\mathcal{X} \to \mathcal{X}'$ が存在する．
\end{lemma}

\begin{proof}
まず Lemma \ref{stacks-limits-lemma-finite-type-closed-in-finite-presentation} の
証明とまったく同じ手順を用い（その記法も借用し），
$\mathcal{X}_i = [U_i/R_i]$ として，
埋め込み $\mathcal{X} \to \mathcal{X}'$ を
射 $\mathcal{X} \to \mathcal{X}' = Y \times_{Y_i} \mathcal{X}_i$
として構成する．
したがって，十分大きい $i$ に対して
$\mathcal{X}_i \to Y_i$ が分離的であることを示せば十分である．
言い換えれば，十分大きい $i$ に対して
$\mathcal{X}_i \to \mathcal{X}_i \times_{Y_i} \mathcal{X}_i$
が固有であることを示せばよい．
射 $U_i \times_{Y_i} U_i \to
\mathcal{X}_i \times_{Y_i} \mathcal{X}_i$ は全射かつ滑らかであり，また
$
R_i = \mathcal{X}_i
\times_{\mathcal{X}_i \times_{Y_i} \mathcal{X}_i} U_i \times_{Y_i} U_i
$
であるから，十分大きい $i$ に対して射
$(s_i, t_i) : R_i \to U_i \times_{Y_i} U_i$
が固有であることを示せば十分である．
Properties of Stacks, Lemma
\ref{stacks-properties-lemma-check-property-covering} を参照されたい．
これを次の段落で証明する．

\medskip\noindent
$U \times_Y U \to Y$ は準分離かつ有限型であることに注意する．
したがって Limits of Spaces, Remark
\ref{spaces-limits-remark-finite-type-gives-well-defined-system}
の構成を用い，ある $i_1 \in I$ と逆系
$(V_i)_{i \geq i_1}$ で
$U \times_Y U = \lim_{i \geq i_1} V_i$ となるものを見つけられる．
Limits of Spaces, Lemma
\ref{spaces-limits-lemma-good-diagram-fibre-product} により，
十分大きい $i$ に対し，射影 $U \times_Y U \to U$ に適用した
構成の関手性から閉埋め込み
$$
V_i \to U_i \times_{Y_i} U_i
$$
が得られる
（厳密には $Y_i$ を $Y \to Y_i$ のスキーム論的像で置き換えるべきなので，
ここには小さな不一致があるが，これは明らかにファイバー積を変えない）．
一方，Limits of Spaces, Lemma
\ref{spaces-limits-lemma-morphism-good-diagram-proper} により，
固有射 $(s, t) : R \to U \times_Y U$
（ここで $\mathcal{X}$ が分離的であることを用いる）に適用した関手性から，
十分大きい $i$ に対して固有な射 $R_i \to V_i$ が得られる．
これらを合成して，すべての十分大きい $i$ に対する固有射
$R_i \to U_i \times_{Y_i} U_i$ を得る．
Limits of Spaces, Remark
\ref{spaces-limits-remark-finite-type-gives-well-defined-system}
の構成の関手性から，十分大きい $i$ に対してこの射は
$(s_i, t_i)$ と同じであることが分かる．これで証明が完了する．
\end{proof}







\section{普遍閉射}
\label{stacks-limits-section-universally-closed}

\noindent
本節は Limits of Spaces, Section
\ref{spaces-limits-section-universally-closed} の類似である．

\begin{lemma}
\label{stacks-limits-lemma-separate}
$g : Z \to Y$ をアフィンスキームの射とする．
$f : \mathcal{X} \to Y$ を代数スタックの準コンパクトな射とする．
$z \in Z$ とし，$T \subset |\mathcal{X} \times_Y Z|$ を
$z \not \in \Im(T \to |Z|)$ を満たす閉部分集合とする．
$\mathcal{X}$ が準コンパクトならば，
$z$ の開近傍 $V \subset Z$，可換図式
$$
\xymatrix{
V \ar[d] \ar[r]_a & Z' \ar[d]^b \\
Z \ar[r]^g & Y,
}
$$
および閉部分集合 $T' \subset |X \times_Y Z'|$ で，
次を満たすものが存在する：
\begin{enumerate}
\item $Z'$ は $Y$ 上有限表示のアフィンスキームである，
\item $z' = a(z)$ とおくと $z' \not \in \Im(T' \to |Z'|)$ である，
\item $|\mathcal{X} \times_Y V|$ における $T$ の逆像は，
$|\mathcal{X} \times_Y V| \to |\mathcal{X} \times_Y Z'|$
によって $T'$ の中に写る．
\end{enumerate}
\end{lemma}

\begin{proof}
これをスキームの射に対する対応する結果から導く．
$\mathcal{X}$ は準コンパクトなので，アフィンスキーム $W$ と
全射な滑らかな射 $W \to \mathcal{X}$ を選べる．
$T_W \subset |W \times_Y Z|$ を $T$ の逆像とする．
このとき $z$ は $T_W$ の像に属さない．
スキームの場合
（Limits, Lemma \ref{limits-lemma-separate}）により，
$z$ の開近傍 $V \subset Z$，スキームの可換図式
$$
\xymatrix{
V \ar[d] \ar[r]_a & Z' \ar[d]^b \\
Z \ar[r]^g & Y,
}
$$
および閉部分集合 $T' \subset |W \times_Y Z'|$ で，
次を満たすものを見つけられる：
\begin{enumerate}
\item $Z'$ は $Y$ 上有限表示のアフィンスキームである，
\item $z' = a(z)$ とおくと $z' \not \in \Im(T' \to |Z'|)$ である，
\item $T_1 = T_W \cap |W \times_Y V|$ は
$|W \times_Y V| \to |W \times_Y Z'|$ によって $T'$ の中に写る．
\end{enumerate}
可換図式
$$
\xymatrix{
W \times_Y Z \ar[d] &
W \times_Y V \ar[l] \ar[rr]_{a_1} \ar[d]_c & &
W \times_Y Z' \ar[d]^q \\
\mathcal{X} \times_Y Z &
\mathcal{X} \times_Y V \ar[l] \ar[rr]^{a_2} & &
\mathcal{X} \times_Y Z'
}
$$
の各正方形はカルテジアンであり，縦の射は全射かつ滑らか，
したがって特に開である．左の正方形を見ると，
$T_1 = T_W \cap |W \times_Y V|$ は，
$T_2 = T \cap |\mathcal{X} \times_Y V|$ の $c$ による逆像である．
Properties of Stacks, Lemma
\ref{stacks-properties-lemma-points-cartesian} により
$a_1(T_1) = q^{-1}(a_2(T_2))$ を得る．
Topology, Lemma
\ref{topology-lemma-open-morphism-quotient-topology} により
$$
q^{-1}\left(\overline{a_2(T_2)}\right) =
\overline{q^{-1}(a_2(T_2))} =
\overline{a_1(T_1)} \subset T'
$$
を得る．$q$ は全射なので，$\overline{a_2(T_2)} \to |Z'|$ の像は
$z'$ を含まない．なぜなら，$T'$ について同じことが成り立つからである．
したがって，上の $Z', V, a, b$ をもつ図式と閉部分集合
$\overline{a_2(T_2)} \subset |\mathcal{X} \times_Y Z'|$ を，
補題が提示する問題の解として取れる．
\end{proof}

\begin{lemma}
\label{stacks-limits-lemma-test-universally-closed}
$f : \mathcal{X} \to \mathcal{Y}$ を代数スタックの準コンパクトな射とする．
次の条件は同値である：
\begin{enumerate}
\item $f$ は普遍閉である，
\item $Z$ がアフィンスキームである局所有限表示の任意の射
$Z \to \mathcal{Y}$ に対し，
写像 $|\mathcal{X} \times_Y Z| \to |Z|$ は閉である，
\item スキーム $V$ と全射な滑らかな射
$V \to \mathcal{Y}$ が存在して，すべての $n \geq 0$ に対し
$|\mathbf{A}^n \times (\mathcal{X} \times_\mathcal{Y} V)|
\to |\mathbf{A}^n \times V|$ は閉である．
\end{enumerate}
\end{lemma}

\begin{proof}
(1) が (2) を含意することは明らかである．

\medskip\noindent
(2) を仮定する．
アフィンスキームの非交和であるスキーム $V$ と，
全射な滑らかな射 $V \to \mathcal{Y}$ を選ぶ．
$f$ が普遍閉であることを示すには，$f$ の基底変換
$\mathcal{X} \times_\mathcal{Y} V \to V$ が普遍閉であることを
示せば十分である．Morphisms of Stacks, Lemma
\ref{stacks-morphisms-lemma-universally-closed-local} を参照されたい．
この基底変換に対しても性質 (2) が成り立つことに注意する．
したがって，(2) が (1) を含意することを証明するためには，
$Y = \mathcal{Y}$ がアフィンスキームであると仮定してよい．

\medskip\noindent
(2) を仮定し，$\mathcal{Y} = Y$ がアフィンスキームであるとする．
$f$ が普遍閉でないならば，
$|\mathcal{X} \times_Y Z| \to |Z|$ が閉でないような
$Y$ 上のアフィンスキーム $Z$ が存在する．
Morphisms of Stacks, Lemma
\ref{stacks-morphisms-lemma-universally-closed-local} を参照されたい．
これは，ある閉部分集合 $T \subset |\mathcal{X} \times_Y Z|$ が存在して，
$\Im(T \to |Z|)$ が閉でないことを意味する．
$T$ の像の閉包に属するが像には属さない $z \in |Z|$ を選ぶ．
Lemma \ref{stacks-limits-lemma-separate} を適用する．
開近傍 $V \subset Z$，可換図式
$$
\xymatrix{
V \ar[d] \ar[r]_a & Z' \ar[d]^b \\
Z \ar[r]^g & Y,
}
$$
および閉部分集合 $T' \subset |\mathcal{X} \times_Y Z'|$ で，
次を満たすものを得る：
\begin{enumerate}
\item $Z'$ は $Y$ 上有限表示のアフィンスキームである，
\item $z' = a(z)$ とおくと $z' \not \in \Im(T' \to |Z'|)$ である，
\item $|\mathcal{X} \times_Y V|$ における $T$ の逆像は，
$|\mathcal{X} \times_Y V| \to |\mathcal{X} \times_Y Z'|$
によって $T'$ の中に写る．
\end{enumerate}
$z'$ は $\Im(T' \to |Z'|)$ の閉包に属すると主張する．
これにより $|\mathcal{X} \times_Y Z'| \to |Z'|$ が閉でないことになり，
(2) の仮定に反する．言い換えれば，この主張から (2) が (1) を含意する．
主張を見るため，次の可換図式を考える：
$$
\xymatrix{
\mathcal{X} \times_Y Z \ar[d] &
\mathcal{X} \times_Y V \ar[l] \ar[d] \ar[r] &
\mathcal{X} \times_Y Z' \ar[d] \\
Z & V \ar[l] \ar[r]^a & Z'
}
$$
$T_V \subset |\mathcal{X} \times_Y V|$ を $T$ の逆像とする．
Properties of Stacks, Lemma
\ref{stacks-properties-lemma-points-cartesian} により，
$|V|$ における $T_V$ の像は，$|Z|$ における $T$ の像の逆像である．
$z$ は $T \to |Z|$ の像の閉包に属し，$|V| \to |Z|$ は開なので，
$z$ は $T_V \to |V|$ の像の閉包に属する．
$|\mathcal{X} \times_Y Z'|$ における $T_V$ の像は $|T'|$ に含まれるから，
$z' = a(z)$ は直ちに $T'$ の像の閉包に属する．

\medskip\noindent
(1) が (3) を含意することは明らかである．
$V \to \mathcal{Y}$ を (3) のとおりとする．
$\mathcal{X} \times_Y V \to V$ が普遍閉であることを示せれば，
Morphisms of Stacks, Lemma
\ref{stacks-morphisms-lemma-universally-closed-local} により
$f$ は普遍閉である．
したがって，$f$ が代数スタックの準コンパクトな射で，
$\mathcal{Y} = Y$ がスキームであり，
すべての $n$ に対して
$|\mathbf{A}^n \times \mathcal{X}| \to |\mathbf{A}^n \times Y|$
が閉であるとき，$f : \mathcal{X} \to \mathcal{Y}$ が
(2) を満たすことを示せば十分である．
$Z$ がアフィンスキームである局所有限表示の射
$Z \to Y$ を取る．
写像 $|\mathcal{X} \times_Y Z| \to |Z|$ が閉であることを示す必要がある．
$Y$ はスキーム，$Z$ はアフィンで，$Z \to Y$ は局所有限表示なので，
埋め込み $Z \to \mathbf{A}^n \times Y$ を見つけられる．
Morphisms, Lemma
\ref{morphisms-lemma-quasi-affine-finite-type-over-S} を参照されたい．
カルテジアン図式
$$
\vcenter{
\xymatrix{
\mathcal{X} \times_Y Z \ar[d] \ar[r] &
\mathbf{A}^n \times \mathcal{X} \ar[d] \\
Z \ar[r] & \mathbf{A}^n \times Y
}
}
\quad
\begin{matrix}
\text{これが誘導する} \\
\text{カルテジアン図式}
\end{matrix}
\quad
\vcenter{
\xymatrix{
|\mathcal{X} \times_Y Z| \ar[d] \ar[r] &
|\mathbf{A}^n \times \mathcal{X}| \ar[d] \\
|Z| \ar[r] & |\mathbf{A}^n \times Y|
}
}
$$
を考える．これは位相空間の図式であり，その横の矢印は局所閉部分集合への
同相写像である
（Properties of Stacks, Lemma
\ref{stacks-properties-lemma-subspace-induced-topology}）．
したがって，$|X \times_Y Z|$ の任意の閉部分集合 $T$ は，
$|\mathbf{A}^n \times Y|$ のある閉部分集合 $T'$ の引き戻しである．
仮定により $|\mathbf{A}^n \times X|$ における $T'$ の像は閉なので，
所望どおり $|Z|$ における $T$ の像は閉である．
\end{proof}
